\chapter{Definitions}
%\addcontentsline{toc}{chapter}{Definitions}

Following terms are used in this work:\par\bigskip

\begin{description}[style=multiline, labelwidth=5cm, leftmargin=5.5cm]

  \item[Multimodal Emotion Recognition]
       The task of identifying human emotions using multiple data modalities such as text, audio, and video.

  \item[Attention Bottleneck Mechanism]
       A fusion mechanism that enables efficient interaction between different modalities through a limited set of shared bottleneck tokens.

  \item[Modality]
       A specific type of input data used in the system, for example textual, acoustic, or visual information.

  \item[Fusion]
       The process of combining information from multiple modalities into a unified representation for emotion prediction.

  \item[Transformer]
       A neural network architecture based on self-attention mechanisms for modeling sequential data and contextual relationships.

  \item[Self-Attention]
       A mechanism that allows a model to determine the importance of different elements within the same sequence.

  \item[Bottleneck Tokens]
       Learnable latent representations used to exchange compressed information between modalities in attention bottleneck models.

  \item[Emotion Classification]
       The process of assigning an input sample to one of several predefined emotion categories.

  \item[Feature Extraction]
       The process of transforming raw input data into informative numerical representations suitable for machine learning models.

  \item[Text Modality]
       Linguistic information represented as textual features or embeddings extracted from speech transcripts.

  \item[Audio Modality]
       Acoustic information such as tone, pitch, energy, and speech characteristics extracted from audio signals.

  \item[Video Modality]
       Visual information including facial expressions, gestures, and movements extracted from video frames.

  \item[BERT]
       A pretrained transformer-based language model used for extracting contextual textual embeddings.

  \item[COVAREP]
       A speech analysis toolkit used for extracting acoustic features from audio data.

  \item[OpenFace]
       An open-source toolkit for facial behavior analysis and extraction of facial action units.

  \item[CMU-MOSEI]
       A large-scale multimodal dataset containing video clips annotated with sentiment and emotion labels.

  \item[Weighted F1-score]
       An evaluation metric that measures classification performance while considering class imbalance.

  \item[Cross-Entropy Loss]
       A loss function commonly used in classification tasks to measure the difference between predicted and true labels.

\end{description}